\documentclass[12pt,a4paper]{article}

\usepackage[utf8]{inputenc}
\usepackage[T1]{fontenc}
\usepackage{lmodern}
\usepackage[margin=1in]{geometry}
\usepackage{setspace}
\usepackage[hidelinks]{hyperref}
\usepackage{amsmath,amssymb}
\usepackage{booktabs}
\usepackage{titlesec}
\usepackage[stretch=15]{microtype}
\usepackage{parskip}
\usepackage[
	backend=biber,
	style=authoryear,
	maxcitenames=3,
	mincitenames=1,
	uniquelist=false,
	sorting=nyt,
	dashed=false,
]{biblatex}
\addbibresource{references.bib}

% Section formatting
\titleformat{\section}{\large\bfseries}{}{0em}{}
\titleformat{\subsection}{\normalsize\bfseries}{}{0em}{}
\titleformat{\subsubsection}{\normalsize\itshape}{}{0em}{}

\title{Technical Note: Certainty Equivalents as a Lens for\\
Characterizing Risk Preference Estimates\\[6pt]
\large A Short Literature Review}
\author{}
\date{March 2026}

\begin{document}
\maketitle
\thispagestyle{empty}

\section{Purpose}

In Section~4.1 of the paper, just before defining the CE mapping, the text
states:

\begin{quote}
``CEs of reference lotteries have been used to assess the post-estimation
reasonableness of RDU and EUT estimates by Harrison et al.\ (2022, p.\
807). They provide a natural way to understand the consequences of
structural models of risk preferences.''
\end{quote}

The single citation is \textcite{Harrison2022}. This note compiles additional
examples---both Harrison and non-Harrison---where certainty equivalents (or
closely related generated quantities such as risk premia) have been used to
characterize, assess, or communicate the implications of structurally estimated
risk preference parameters.

The uses fall into several categories: (i)~translating estimated structural
parameters into CEs or risk premia to make results interpretable;
(ii)~computing CEs from estimated preferences to perform welfare analysis;
(iii)~using prior predictive CEs to calibrate or validate prior
specifications; and (iv)~comparing CEs implied by alternative structural
models (EUT vs.\ RDU vs.\ CPT) to evaluate model adequacy.

\section{Welfare Evaluation via Certainty Equivalents}

The most developed use of CEs as a post-estimation diagnostic is in behavioral
welfare economics, where the CE of a risky prospect---computed from an
individual's estimated preference parameters---provides a dollar-denominated
measure of welfare that can be compared across individuals, models, and policy
alternatives.

\textcite{Harrison2016} introduced this approach explicitly. They estimate
individual risk preferences (CRRA + Prelec PWF under RDU) from one
experimental task and then compute the CE of risky insurance policies faced by
each individual. Comparing this CE to the insurance premium yields an
individual-specific measure of expected consumer surplus. Their central finding
is that the metric of insurance take-up, which is widely used in policy
evaluations, provides a qualitatively incorrect guide to welfare effects---a
conclusion that is only visible when results are translated from parameter
space to CE space.

\textcite{Gao2022} extend this framework to a Bayesian hierarchical model.
They compute the CE of insurance options for each individual from posterior
draws of that individual's risk preference parameters, yielding a full
posterior distribution over consumer surplus. The hierarchical structure
provides informative priors for individuals with limited choice data, and the
welfare conclusions are materially different from those obtained with diffuse
priors. Again, the CE is the object through which estimated risk preferences
are given economic meaning.

\textcite{Barseghyan2013} use administrative data on insurance deductible
choices to separately identify utility curvature and probability weighting
within a structural model. While they do not use the Prelec parameterization
specifically, their identification strategy exploits the fact that the CE
implications of CRRA curvature and probability weighting differ across
deductible levels. Implied risk premia from the estimated model are used to
assess which component---diminishing marginal utility or probability
distortion---drives observed risk aversion in the field.

\section{CEs to Interpret and Communicate Structural Estimates}

A second class of uses treats CEs as a more interpretable summary of what
structural parameter estimates actually mean. Estimated values of $r$, $\phi$,
and $\eta$ are difficult to evaluate in isolation---their economic content is
opaque without translation into quantities with natural units. CEs provide that
translation.

\textcite{HarrisonRutstrom2008}, in their comprehensive methodological survey,
use implied CEs and risk premia to illustrate what different CRRA coefficient
estimates mean in practice. They compute the CE of specific reference lotteries
at various values of $r$ to give readers a concrete sense of what ``$r = 0.67$''
implies about economic behavior. This pedagogical use of CEs has been widely
adopted in subsequent experimental economics research.

\textcite{HarrisonSwarthout2023} compare EUT, RDU, Dual Theory, and CPT by
computing the CE of each subject's choices under each estimated model. The
comparison of model-implied CEs is used to assess which model best
characterizes individual behavior: RDU-implied CEs are more consistent with
observed choices than CPT-implied CEs, leading to their conclusion that RDU is
the best single characterization of individual choice.

\textcite{LauYoo2025}, in a recent \textit{Econometrica} paper on higher-order
risk preferences, compute predicted risk premia under both EUT and RDU for
higher-order risk apportionment tasks. They find that predicted risk premia
under EUT and RDU gradually converge as the order of risk increases, with RDU
predicting substantially larger first-order risk premia (mean 7.2\% vs.\ 1.4\%
under EUT). The risk premium---which is the expected value minus the CE---is
the vehicle through which these structural estimates are given economic content.

\textcite{Andersen2008} use the structural model to compute the implied
discount rate that is consistent with both estimated risk preferences and
observed intertemporal choices. While the generated quantity here is a discount
rate rather than a CE, the logic is identical: structural parameter estimates
are translated into an economically interpretable quantity to assess whether
the estimates are reasonable and to compare across specifications.

\section{Prior Predictive CEs for Calibration}

A third use treats CEs as a diagnostic for prior specification rather than
posterior inference. Here, the researcher computes the CE implied by draws from
the prior distribution and checks whether the prior predictive distribution of
CEs is reasonable.

\textcite{Bland2025} advocates this approach explicitly, proposing that priors
for structural models in economic experiments should be calibrated by examining
their implications for certainty equivalents. In a Monte Carlo exercise, he
shows that a Bayesian estimator with a CE-calibrated prior estimates both
structural parameters and certainty equivalents with smaller error than maximum
likelihood or estimators with diffuse priors. The CE is used both as a
calibration target (pre-estimation) and as an evaluation metric
(post-estimation).

The present paper extends this idea to the hierarchical setting: the prior
predictive CE distribution is used to verify that the hyperprior specification
implies reasonable beliefs about individual behavior, and the posterior
predictive CE distribution is used to characterize what the model has learned
from the data.

\section{CEs in Non-Parametric and Semi-Parametric Elicitation}

\textcite{Abdellaoui2000} developed a parameter-free method for eliciting
utility and probability weighting functions under RDU. The procedure is
grounded in certainty equivalents: subjects state the certain amount that
makes them indifferent to a lottery, and these CEs are used to trace out the
utility function and weighting function nonparametrically. This is
conceptually the inverse of the post-estimation CE computation---rather than
going from estimated parameters to implied CEs, one goes from elicited CEs to
inferred preference components---but it rests on the same foundational
insight: CEs are the natural observable through which risk preferences are
revealed and characterized.

\section{CEs Across Populations and Contexts}

\textcite{BruhinFehrDudaEpper2010} estimate a finite mixture model for
probability weighting using data from Swiss and Chinese experiments. They
classify subjects into types based on the degree of probability distortion and
then characterize the types partly through their implied CEs for reference
lotteries. The CE provides a common metric for comparing what ``type 1'' and
``type 2'' behavior means in economic terms, even though the types are defined
in the parameter space.

\textcite{Tanaka2010} estimate structural risk preferences for Vietnamese
farmers and translate the estimated parameters into risk premia (equivalently,
deviations of CEs from expected values) to assess the economic significance of
the estimates and to link laboratory-estimated preferences to household survey
data on real economic decisions.

\textcite{vonGaudecker2011} structurally estimate utility curvature, loss
aversion, and resolution timing preferences in a representative Dutch sample.
They characterize heterogeneity partly through predicted choice probabilities
and implied risk attitudes, using the structural estimates to compute
quantities that reveal what the estimated parameter heterogeneity means for
economic behavior.

\section{Summary}

The use of CEs to characterize risk preference estimates is pervasive, though
it is rarely discussed as a methodological principle in its own right. The
examples above can be grouped as follows:

\begin{tabular}{p{4cm} p{5cm} p{5cm}}
\toprule
\textbf{Use} & \textbf{Examples} & \textbf{Role of CE} \\
\midrule
Welfare evaluation & \textcite{Harrison2016}; \textcite{Gao2022};
  \textcite{Barseghyan2013} &
  Dollar-valued consumer surplus from estimated preferences \\[6pt]
Interpreting estimates & \textcite{HarrisonRutstrom2008};
  \textcite{HarrisonSwarthout2023}; \textcite{LauYoo2025} &
  Translate abstract parameters into economic quantities \\[6pt]
Model comparison & \textcite{HarrisonSwarthout2023}; \textcite{LauYoo2025} &
  Compare CEs/risk premia implied by EUT vs.\ RDU \\[6pt]
Prior calibration & \textcite{Bland2025} &
  Verify prior reasonableness via prior predictive CE \\[6pt]
Nonparametric elicitation & \textcite{Abdellaoui2000} &
  Infer preferences from observed CEs \\[6pt]
Cross-population comparison & \textcite{BruhinFehrDudaEpper2010};
  \textcite{Tanaka2010}; \textcite{vonGaudecker2011} &
  Common metric for comparing types/populations \\
\bottomrule
\end{tabular}

\vspace{1em}

The paper currently cites only \textcite{Harrison2022} for this claim. The
examples above---spanning welfare economics, model comparison, prior
calibration, and nonparametric elicitation---demonstrate that the practice is
far more widespread. I would suggest expanding the sentence to cite several of
these, especially those from non-Harrison research groups, to convey the
breadth of the practice. A possible revision:

\begin{quote}
``CEs of reference lotteries have been used to assess the post-estimation
reasonableness of RDU and EUT estimates
\parencite[807]{Harrison2022}, to evaluate the welfare
implications of estimated risk preferences
\parencite{Harrison2016, Gao2022, Barseghyan2013}, to compare the
economic content of alternative structural models
\parencite{HarrisonSwarthout2023, LauYoo2025}, and to calibrate
prior specifications \parencite{Bland2025}. They provide a natural way to
understand the consequences of structural models of risk preferences.''
\end{quote}

\printbibliography

\end{document}
